\documentclass[UTF8,fontset=fandol,11pt]{ctexart}
\usepackage[a4paper,margin=24mm]{geometry}
\usepackage{amsmath,amssymb,amsthm,graphicx,tabularx,booktabs,hyperref}
\newtheorem{definition}{定义}[section]
\newtheorem{theorem}[definition]{定理}
\newtheorem{example}[definition]{例题}
\newtheorem{exercise}[definition]{练习}
\hypersetup{unicode=true,colorlinks=true,linkcolor=blue}
\title{Offline draft recovery}
\date{}
\begin{document}
\maketitle
\par\noindent 离线合成演示：不代表真实模型对本专题的教学验收。\par
\section{Definition}\label{sec:S01}
% UNIT S01.U01
\begin{definition}
最小二乘法以残差平方和作为待最小化的目标函数。
\end{definition}
% UNIT S01.U04
\begin{proof}
本证明需要将给定候选参数代入两条正规方程验证。 已完成独立修复并逐项核验。
\end{proof}
% UNIT S01.U10
\par\noindent\textbf{解：}本练习答案的有限求和及相应分数结果需要重新核对。 已完成独立修复并逐项核验。\par
\end{document}
